% Corrigé intégré automatiquement pour la banque Exercices.
% Source : A_Integrer/DDS_04/076_Geometrie/076_Geometrie.tex
\begin{corrigebox}[Corrigé]
\CorrigeQuestion{1}
$x(t)=\ell_1 \cos\theta +\ell_2 \sqrt{1-r^2\sin^2\theta}$.
\CorrigeQuestion{2}
$A=\sqrt{1-r^2\sin^2\theta}$, $B=-\ell_1\sin\theta\sqrt{1-r^2\in^2\theta}-\ell_2r^2\sin\theta\cos\theta$.
\CorrigeQuestion{3}
$\theta=0\degres$.
\CorrigeQuestion{4}
La résolution consiste à identifier les données et l'inconnue, choisir la loi du cours adaptée, écrire l'équation symbolique avant toute application numérique, puis vérifier l'homogénéité, le signe et l'ordre de grandeur du résultat. Les hypothèses utilisées doivent être explicitement contrôlées à la fin.
\CorrigeQuestion{5}
La résolution consiste à identifier les données et l'inconnue, choisir la loi du cours adaptée, écrire l'équation symbolique avant toute application numérique, puis vérifier l'homogénéité, le signe et l'ordre de grandeur du résultat. Les hypothèses utilisées doivent être explicitement contrôlées à la fin.
\CorrigeQuestion{6}
La résolution consiste à identifier les données et l'inconnue, choisir la loi du cours adaptée, écrire l'équation symbolique avant toute application numérique, puis vérifier l'homogénéité, le signe et l'ordre de grandeur du résultat. Les hypothèses utilisées doivent être explicitement contrôlées à la fin.
\end{corrigebox}
